% Paper 26 -- Section 7: Evaluation
%
% Live measurements from artifacts in research/paper-26-artifacts/:
%   tropical-gamma-experiment-2026-05-24.json  -- §7.1 (C1 tropical γ)
%   jepa-training-experiment-2026-05-24.json   -- §7.2 (C2 JEPA convergence)
%   grpo-diversity-experiment-2026-05-24.json  -- §7.3 (C3 GRPO diversity)
%   c1-token-reduction-2026-05-24.json         -- §7.4 (token routing)
%   section7-harness-2026-05-21.json           -- §7.5–7.8 (substrate benchmarks)
%
% To regenerate substrate benchmarks:
%   pnpm run bench:paper26:section7

\section{Evaluation}
\label{sec:evaluation}

We evaluate the paper's three primary contributions first
(\S\,\ref{sec:eval-tropical-gamma}--\ref{sec:eval-token-routing}),
then benchmark the sovereign code-intelligence substrate that generates
training slices for those contributions
(\S\,\ref{sec:eval-holograph}--\ref{sec:eval-holomesh-coordination}).
All experiments use unmodified production code; no reimplementation or
harness substitution.

% ─────────────────────────────────────────────────────────────────────────────
\subsection{Tropical-Geometry Coordination Agreement (Contribution 1)}
\label{sec:eval-tropical-gamma}
% ─────────────────────────────────────────────────────────────────────────────

\paragraph{Setup.}
Two real \texttt{physics}-domain hemisphere Pillars (via
\texttt{makeParallelPillar()} from \texttt{ParallelPillar.ts}) observe the
same true state $(s_1, s_2) \in [0,1]^2$ with zero-mean Gaussian measurement
noise of std $= d$, split into shared common-mode ($\rho = 0.5$) and
independent parts.  Each $(left, right)$ pair is fed through
\texttt{generateParallel()};  $\gamma = 1 - \mathrm{box\_area}$ is
cross-checked against the standalone \texttt{computeParallelBounds} export
(assertion: identical values).  Sweep: 11 divergence levels $d \in [0, 1.0]$,
1,200 samples each, deterministic (mulberry32, seed 262026); $N = 13{,}200$
total.

\paragraph{Claim under test.}
$\gamma$ is a coordination-agreement metric: $\gamma = 1$ when the two
hemispheres observe the same state, and \emph{monotonically decreasing} as
their measurements diverge.

\paragraph{Results -- Table~\ref{tab:tropical-gamma}.}

\begin{table}[h]
\centering
\caption{Tropical bilateral $\gamma$ vs measurement divergence ($N = 13{,}200$,
  deterministic, seed 262026).  $\gamma = 1.000$ at zero divergence (degenerate
  bounding box); Pearson $r(\gamma, d) = -0.348$ confirms the
  coordination-agreement property.
  Artifact: \texttt{tropical-gamma-experiment-2026-05-24.json}.}
\label{tab:tropical-gamma}
\begin{tabular}{rrr}
\toprule
Divergence $d$ & Mean $\gamma$ & Mean box area \\
\midrule
0.00 & 1.000 & 0.000 \\
0.20 & 0.981 & 0.019 \\
0.50 & 0.936 & 0.064 \\
1.00 & 0.915 & 0.085 \\
\bottomrule
\end{tabular}
\end{table}

\noindent
$\gamma$ is monotonically decreasing across all 11 divergence levels
(all-steps monotone check: \textbf{pass}).  $\gamma = 1.000$ at $d = 0$
and $\gamma = 0.915$ at $d = 1.0$; Pearson $r = -0.348$ across all samples.
The coordination-agreement claim is validated.

% ─────────────────────────────────────────────────────────────────────────────
\subsection{PillarJEPA Convergence (Contribution 2)}
\label{sec:eval-jepa-convergence}
% ─────────────────────────────────────────────────────────────────────────────

\paragraph{Setup.}
A standalone gradient-descent experiment drives the real \texttt{JEPAPredictor}
MLP (from \texttt{JEPAPredictor.ts}) with exact analytic backprop and Adam
($\eta = 0.001$).  Gradients are verified against central finite differences
(rel-err $= 0.000 < 10^{-3}$, confirming the backprop is exact).  Training
runs for 1,500 steps with minibatch size 16; evaluation uses a fixed 256-sample
held-out batch evaluated every 25 steps (61 eval points total).

\paragraph{Claim under test.}
\textsc{PillarJEPA} trains stably: held-out eval loss decreases, no NaN or Inf,
tail plateau is stable, temporal gating relieves conservation pressure as
$\kappa \to 1$.

\paragraph{Results.}
Held-out eval loss decreases from $0.3406$ to $0.1752$ ($1.944\times$
reduction, $48.6\%$; artifact field: \texttt{eval\_reduction\_factor}).
Stability checks:

\begin{itemize}
  \item No NaN or Inf at any step: \textbf{pass}
  \item Held-out eval reduction $\geq 30\%$ (actual $48.6\%$): \textbf{pass}
  \item Tail plateau CV $= 4\%$ ($< 25\%$ threshold, no upward drift): \textbf{pass}
  \item Smoothed eval curve non-increasing: \textbf{pass}
  \item Temporal gating: $\lambda_c^{\mathrm{eff}}$ decreases from $0.100$ to
        $0.000$ as $\kappa \to 1$, consistent with Eq.~\ref{eq:lambda-eff}: \textbf{pass}
\end{itemize}

All 5 stability checks pass.
Artifact: \texttt{jepa-training-experiment-2026-05-24.json}.

\paragraph{Scope boundary.}
The trained predictor weights are validated via the real
\texttt{pillarJepaHandler}: the handler emits a finite, well-formed
\texttt{pillarjepa:loss} with all components present, confirming the trained
predictor is a valid drop-in for the production trait.  The handler uses its
own context encoder (EmbeddingTrait hash), so its MSE term differs from the
training loop's by construction (expected encoder gap, max delta $0.172$);
this is not a discrepancy in the objective.

% ─────────────────────────────────────────────────────────────────────────────
\subsection{GRPO Slice Diversity (Contribution 3)}
\label{sec:eval-grpo-diversity}
% ─────────────────────────────────────────────────────────────────────────────

\paragraph{Setup.}
The real \texttt{sliceEmitterHandler} (from \texttt{SliceEmitter.ts}) and the
real seed Pillars (\texttt{PillarRegistry.ts}) are driven over two streams of
$W = 1{,}000$ slices: a healthy diverse stream and a mode-collapsed stream (all
slices identical).  Fingerprint: \texttt{axis\_1:axis\_2:p1.toFixed(2):p2.toFixed(2)}.

\paragraph{Claim under test.}
Diverse Pillar coordinates structurally guarantee diverse policy gradients;
mode collapse is detectable and carries near-zero gradient signal.

\paragraph{Results.}
The cross-stream diversity discrimination ratio at $N = 1{,}000$ is
$371\times$ (371 unique fingerprints vs 1).  The cross-stream reward-variance
ratio is $2.29 \times 10^{5}$ (healthy $4.15 \times 10^{-2}$
vs collapsed $1.81 \times 10^{-7}$), confirming that diverse streams carry a
non-trivial reward signal while collapsed streams carry $\approx 0$.
Artifact: \texttt{grpo-diversity-experiment-2026-05-24.json}.

\paragraph{Honest finding.}
The static alert threshold ($\rho > 0.8$) in \textsc{SliceEmitter} is
miscalibrated for the lifetime-cumulative fingerprint ratio: the healthy
agent's $\rho$ also falls to $0.371$ at $N = 1{,}000$, triggering a false
positive.  The robust collapse signal is the discrimination ratio and the
unique-count growth rate; the recommended fix (rolling-window $\rho$ over
the last $W$ slices) is documented in Appendix~\ref{app:grpo}.

% ─────────────────────────────────────────────────────────────────────────────
\subsection{Token Routing Efficiency}
\label{sec:eval-token-routing}
% ─────────────────────────────────────────────────────────────────────────────

\paragraph{Setup.}
Total coordination-token counts are measured for the Pillar-Slice board-delta
protocol vs a na\"{i}ve full-transcript rebroadcast baseline, using real BPE
tokenisation (\texttt{cl100k\_base}, local, no external API).  The na\"{i}ve
baseline rebroadcasts the full accumulated transcript to all $N - 1$ peers each
round ($O(N^2 H)$ total tokens); the Pillar-Slice baseline emits one
\texttt{SemanticCollaborationMessage} delta per agent per round to a shared
board read once ($O(NH)$ total tokens).  Both convey the same decision;
the baseline is not inflated.
Artifact: \texttt{c1-token-reduction-2026-05-24.json}
(\texttt{canonical\_operating\_point} field).

\paragraph{Results.}
At the canonical operating point ($N = 100$ agents, $H = 50$ rounds):
slice-protocol total $678{,}642$ tokens; broadcast total $810{,}836{,}334$
tokens; reduction $99.9\%$.  The reduction is a \emph{routing} effect
(fan-out avoidance and no history resend), not per-message encoding
compression.  At the per-message level, the compact latent tuple ($21$ tokens)
is $67.2\%$ smaller than a natural-language utterance ($64$ tokens) conveying
the same decision; the full JSON envelope ($136$ tokens) is $112\%$
\emph{larger} than NL.  The asymptotic advantage is $O(NH)$ vs $O(N^2H)$;
the compact-tuple encoding reduction is a secondary benefit.

\paragraph{Comparison with prior work.}
\label{sec:eval-latentmas-comparison}
LatentMAS~\cite{latentmas2025} (Hu et al., ICML 2026 Spotlight) achieves
83.7\% token reduction on general natural-language multi-agent tasks;
RecursiveMAS~\cite{recursivemas2026} independently validates latent
inter-agent coordination on similar tasks.  \PillarSlice{} does not offer a
competing result on those task classes: every coordination event is gated on a
live physics-solver output and anchored by a cryptographic
\texttt{SimulationContract} receipt, so the figures are not directly
comparable.  \PillarSlice{} offers a complementary capability: every
coordination event is post-hoc verifiable from solver state to decision via
cryptographic receipt replay---a guarantee neither system provides and that is
necessary for safety-critical simulation pipelines.  HoloScript is fully
open-source under a permissive license; neither LatentMAS nor RecursiveMAS
releases a production runtime.

Our architecture predates the RecursiveMAS submission: the open-source
\texttt{uaa2-service} repository (first commit 2025-10-30) predates the
2026-04-28 arXiv submission by six months, providing independent confirmation
that latent inter-agent coordination is a convergent design.

% ─────────────────────────────────────────────────────────────────────────────
\subsection{HoloGraph: O(1) Event-Chain Lookup}
\label{sec:eval-holograph}
% ─────────────────────────────────────────────────────────────────────────────

HoloGraph provides the O(1) event-chain lookup that enables training-slice
generation at production scale.  The following benchmark validates this
substrate claim.

\paragraph{Setup.}
We generate synthetic \texttt{CodebaseGraph} instances via \texttt{SyntheticGraphFactory}
with controlled event topology and known ground truth (\texttt{EventGroundTruth}).
Each graph has $N$ source files, $E$ distinct event names, and 3 listener files
per event (1 emitter + 3 listeners = 4 ground-truth files per event chain).
We compare two retrieval paths: \textbf{HoloGraph}
(\texttt{getEventChain(eventName)}, O(1) hash-map lookup) vs
\textbf{Embedding scan} (\texttt{EmbeddingIndex.search(query,~10)},
O($N \cdot D$) cosine similarity).
Latency: median over 500 iterations (per-event, warm JIT).

\paragraph{Results -- Table~\ref{tab:holograph-lookup}.}

\begin{table}[h]
\centering
\caption{Event-chain lookup: HoloGraph O(1) vs Embedding O($N \cdot D$).
  Embedding baseline: \textsc{StructuralEmbeddingProvider} (384-dim, FNV-1a hash,
  no semantic training, no API key).  HoloGraph recall is exact by construction.
  Artifact: \texttt{section7-harness-2026-05-21.json}.}
\label{tab:holograph-lookup}
\begin{tabular}{rrrrrrr}
\toprule
Files & Syms & Events & HG $\mu$s & Emb $\mu$s & HG recall & Emb recall@10 \\
\midrule
   50 &  200 &     10 &  0.950 &   302.3 & 1.000 & 0.125 \\
  500 & 2000 &     50 &  4.145 &  2365.1 & 1.000 & 0.025 \\
 2000 & 8000 &    100 &  7.960 & 12416.8 & 1.000 & 0.025 \\
\bottomrule
\end{tabular}
\end{table}

HoloGraph achieves 100\% recall at sub-microsecond to single-digit-microsecond
latency.  Embedding scan recall degrades from 12.5\% to 2.5\% as graph size
grows, because the query \texttt{"handler for event X"} has no semantic link
to the generic symbol names (\texttt{fn0}, \texttt{fn1}, \ldots) in the synthetic
corpus.  The speedup ratio grows from 318$\times$ (50 files) to 571$\times$
(500 files) to 1560$\times$ (2000 files), consistent with O($N \cdot D$)
scaling.

\paragraph{Ablation: semantic model vs structural baseline.}
\label{sec:ablation-xenova}

A reviewer might argue that \textsc{StructuralEmbeddingProvider} is a weak
baseline because it uses no semantic training.  We address this with a Xenova
\texttt{all-MiniLM-L6-v2} ablation (384-dim, trained on natural language,
no API key).

\begin{table}[h]
\centering
\caption{Ablation: Xenova semantic model vs StructuralEmbeddingProvider on the
  50-file synthetic corpus.  Both providers achieve identical recall (12.5\%),
  because the corpus symbol names carry no semantic signal about event
  membership.  The ablation isolates the latency claim.}
\label{tab:ablation-xenova}
\begin{tabular}{lrrrrrr}
\toprule
Provider & HG $\mu$s & Emb $\mu$s & Speedup & HG recall & Emb recall@10 \\
\midrule
StructuralEmbeddingProvider &  0.950 &   302.3 &   318$\times$ & 1.000 & 0.125 \\
Xenova/all-MiniLM-L6-v2    &  1.650 &  3880.8 &  2352$\times$ & 1.000 & 0.125 \\
\bottomrule
\end{tabular}
\end{table}

Both embedding providers achieve identical recall (12.5\%): the synthetic symbols
(\texttt{fn0}--\texttt{fn3}) have no lexical connection to event names.
The ablation establishes one claim only: \emph{the O(1) latency advantage holds
regardless of embedding quality}.  Xenova's inference time makes it $13\times$
slower than the structural embedding baseline and $2352\times$ slower than
HoloGraph on the same 50-file corpus.

For NL$\to$code recall on semantically meaningful symbol names, see
\S\ref{sec:eval-holoembed}.

% ─────────────────────────────────────────────────────────────────────────────
\subsection{HoloEmbed: NL$\to$Code Recall Without an External API}
\label{sec:eval-holoembed}
% ─────────────────────────────────────────────────────────────────────────────

\paragraph{Setup.}
We construct a 50-symbol recall corpus: 10 \emph{target} symbols with
semantically meaningful PascalCase names (e.g., \texttt{PillarSliceEmitter},
\texttt{BrainCoordNodeMapper}) and 40 \emph{distractor} symbols from an
unrelated WebGL domain.  For each target, the NL query is the
\texttt{camelSplit} of its name.  Recall@10 measures how many targets appear
in the top-10 results from the 50-symbol indexed corpus.

\paragraph{Results -- Table~\ref{tab:holoembed-recall}.}

\begin{table}[h]
\centering
\caption{NL$\to$code recall@10 on a 50-symbol corpus (10 named targets + 40
  WebGL distractors).  HoloEmbed uses 768-dim char-trigram subword features;
  Structural uses 384-dim topology features only.}
\label{tab:holoembed-recall}
\begin{tabular}{llr}
\toprule
Provider & Architecture & Recall@10 \\
\midrule
StructuralEmbeddingProvider & 384-dim topology (FNV-1a, no name encoding) & 0.0\% \\
HoloEmbedProvider           & 768-dim structural + char-trigram subwords  & 90.0\% \\
\bottomrule
\end{tabular}
\end{table}

\noindent
The structural baseline achieves 0\% recall: its topology features carry no
lexical information, so query tokens find no alignment.  The random baseline
for recall@10 over 50 symbols is 20\%; structural recall at 0\% is below
random, confirming active misalignment with NL queries.

HoloEmbedProvider achieves 90\% recall (9 of 10 targets retrieved in top-10)
via three 128-bin char-trigram blocks over symbol name, docComment, and event
names.  Overlapping 3-character windows allow NL tokens such as
\texttt{"pillar"} to match \texttt{"PillarSliceEmitter"} via shared trigrams
(\texttt{pil}, \texttt{ill}, \texttt{lla}, \ldots).  The 90-percentage-point
improvement is achieved with zero external dependencies, zero API keys, and
zero model downloads.  We do not report a Xenova comparison on this task; the
embedding-latency tradeoffs for semantic models are captured in the
event-chain ablation (Table~\ref{tab:ablation-xenova}).

% ─────────────────────────────────────────────────────────────────────────────
\subsection{Real-Corpus Generalization}
\label{sec:eval-real-corpus}
% ─────────────────────────────────────────────────────────────────────────────

We scan two real TypeScript corpora from the HoloScript workspace using
\texttt{CodebaseScanner} and \texttt{CodebaseGraph.buildFromScanResult()}: Corpus A
(\texttt{absorb-service/src/engine/}, $\approx$20 files) and Corpus B
(\texttt{absorb-service/src/}, $\approx$60 files), both with no synthetic fixtures.

On both corpora, HoloGraph query latency is strictly less than embedding scan
latency, consistent with the synthetic benchmark.  Structural integrity holds
for all discovered EventEdges.  Both corpora contain zero
\texttt{.emit()} / \texttt{.on()} pattern pairs in the scanned scope (the
codebase uses a different event system); an O(1) hash-map miss is still faster
than an O($N \cdot D$) embedding scan, regardless of graph sparsity.

% ─────────────────────────────────────────────────────────────────────────────
\subsection{HoloMesh Coordination Latency}
\label{sec:eval-holomesh-coordination}
% ─────────────────────────────────────────────────────────────────────────────

A read-only HoloMesh team-protocol probe measures the live board endpoint used
by room agents (\texttt{GET /api/holomesh/team/:teamId/board}).

\begin{table}[h]
\centering
\caption{HoloMesh team-protocol coordination latency.  Sequential path: 8
  authenticated board reads; burst path: 6 concurrent board reads.
  Artifact: \texttt{section7-harness-2026-05-21.json}.}
\label{tab:holomesh-coordination}
\begin{tabular}{lrrrrr}
\toprule
Mode & p50 ms & p95 ms & seq req/s & burst req/s & success \\
\midrule
live-holomesh & 53.159 & 222.387 & 12.039 & 19.740 & 100.0\% \\
\bottomrule
\end{tabular}
\end{table}

% ─────────────────────────────────────────────────────────────────────────────
\subsection{Threats to Validity and External Validation}
\label{sec:eval-threats}
% ─────────────────────────────────────────────────────────────────────────────

\paragraph{Internal validity.}
Synthetic corpora (\S\ref{sec:eval-holograph}) use controlled topologies with
known ground truth, which may not reflect event-chain distributions in
unfamiliar codebases.  Real-corpus generalization (\S\ref{sec:eval-real-corpus})
partially addresses this.  The contribution-validation experiments
(\S\ref{sec:eval-tropical-gamma}--\ref{sec:eval-grpo-diversity}) use
deterministic seeds and fixed held-out batches; variability across random seeds
and larger corpora is future work.

\paragraph{External validation of the filesystem-as-ground-truth design.}
The core design choice underlying \HoloGraph\
--- that plain, version-controlled files serve as the agent's
authoritative context store rather than opaque in-memory state or
framework-managed DAGs --- has received independent practitioner
validation.  Van Clief's Interpretable Context Methodology
(ICM)~\cite{vanclief2026icm}, which applies the same principle to
workflow pipelines, reports adoption by over 7,000 practitioners in
its first 12 days of community launch~\cite{cliefnotes2026}.  ICM's
numbered-folder stage contracts (\texttt{01-research/},
\texttt{02-script/}, \dots) and \texttt{CONTEXT.md} interface tables
mirror HoloScript's \texttt{CLAUDE.md} + \texttt{skills/} +
\texttt{MEMORY.md} + hook architecture, confirming that
filesystem-grounded agent coordination is not an isolated design
choice but a convergent practice across independent practitioner
communities.
